% Options for packages loaded elsewhere
\PassOptionsToPackage{unicode}{hyperref}
\PassOptionsToPackage{hyphens}{url}
\PassOptionsToPackage{table}{xcolor}
%
\documentclass[
  letterpaper,
]{article}
\usepackage{amsmath,amssymb}
\usepackage{iftex}
\usepackage[T1]{fontenc}
\usepackage[utf8]{inputenc}
\usepackage{textcomp}
\usepackage{lmodern}
\ifPDFTeX\else
  % xetex/luatex font selection
\fi
% Use upquote if available, for straight quotes in verbatim environments
\IfFileExists{upquote.sty}{\usepackage{upquote}}{}
\IfFileExists{microtype.sty}{% use microtype if available
  \usepackage[]{microtype}
  \UseMicrotypeSet[protrusion]{basicmath} % disable protrusion for tt fonts
}{}
\makeatletter
\@ifundefined{KOMAClassName}{% if non-KOMA class
  \IfFileExists{parskip.sty}{%
    \usepackage{parskip}
  }{% else
    \setlength{\parindent}{0pt}
    \setlength{\parskip}{6pt plus 2pt minus 1pt}}
}{% if KOMA class
  \KOMAoptions{parskip=half}}
\makeatother
\usepackage{xcolor}
\usepackage{color}
\usepackage{fancyvrb}
\newcommand{\VerbBar}{|}
\newcommand{\VERB}{\Verb[commandchars=\\\{\}]}
\DefineVerbatimEnvironment{Highlighting}{Verbatim}{commandchars=\\\{\}}
% Add ',fontsize=\small' for more characters per line
\newenvironment{Shaded}{}{}
\newcommand{\AlertTok}[1]{\textcolor[rgb]{1.00,0.00,0.00}{\textbf{#1}}}
\newcommand{\AnnotationTok}[1]{\textcolor[rgb]{0.38,0.63,0.69}{\textbf{\textit{#1}}}}
\newcommand{\AttributeTok}[1]{\textcolor[rgb]{0.49,0.56,0.16}{#1}}
\newcommand{\BaseNTok}[1]{\textcolor[rgb]{0.25,0.63,0.44}{#1}}
\newcommand{\BuiltInTok}[1]{\textcolor[rgb]{0.00,0.50,0.00}{#1}}
\newcommand{\CharTok}[1]{\textcolor[rgb]{0.25,0.44,0.63}{#1}}
\newcommand{\CommentTok}[1]{\textcolor[rgb]{0.38,0.63,0.69}{\textit{#1}}}
\newcommand{\CommentVarTok}[1]{\textcolor[rgb]{0.38,0.63,0.69}{\textbf{\textit{#1}}}}
\newcommand{\ConstantTok}[1]{\textcolor[rgb]{0.53,0.00,0.00}{#1}}
\newcommand{\ControlFlowTok}[1]{\textcolor[rgb]{0.00,0.44,0.13}{\textbf{#1}}}
\newcommand{\DataTypeTok}[1]{\textcolor[rgb]{0.56,0.13,0.00}{#1}}
\newcommand{\DecValTok}[1]{\textcolor[rgb]{0.25,0.63,0.44}{#1}}
\newcommand{\DocumentationTok}[1]{\textcolor[rgb]{0.73,0.13,0.13}{\textit{#1}}}
\newcommand{\ErrorTok}[1]{\textcolor[rgb]{1.00,0.00,0.00}{\textbf{#1}}}
\newcommand{\ExtensionTok}[1]{#1}
\newcommand{\FloatTok}[1]{\textcolor[rgb]{0.25,0.63,0.44}{#1}}
\newcommand{\FunctionTok}[1]{\textcolor[rgb]{0.02,0.16,0.49}{#1}}
\newcommand{\ImportTok}[1]{\textcolor[rgb]{0.00,0.50,0.00}{\textbf{#1}}}
\newcommand{\InformationTok}[1]{\textcolor[rgb]{0.38,0.63,0.69}{\textbf{\textit{#1}}}}
\newcommand{\KeywordTok}[1]{\textcolor[rgb]{0.00,0.44,0.13}{\textbf{#1}}}
\newcommand{\NormalTok}[1]{#1}
\newcommand{\OperatorTok}[1]{\textcolor[rgb]{0.40,0.40,0.40}{#1}}
\newcommand{\OtherTok}[1]{\textcolor[rgb]{0.00,0.44,0.13}{#1}}
\newcommand{\PreprocessorTok}[1]{\textcolor[rgb]{0.74,0.48,0.00}{#1}}
\newcommand{\RegionMarkerTok}[1]{#1}
\newcommand{\SpecialCharTok}[1]{\textcolor[rgb]{0.25,0.44,0.63}{#1}}
\newcommand{\SpecialStringTok}[1]{\textcolor[rgb]{0.73,0.40,0.53}{#1}}
\newcommand{\StringTok}[1]{\textcolor[rgb]{0.25,0.44,0.63}{#1}}
\newcommand{\VariableTok}[1]{\textcolor[rgb]{0.10,0.09,0.49}{#1}}
\newcommand{\VerbatimStringTok}[1]{\textcolor[rgb]{0.25,0.44,0.63}{#1}}
\newcommand{\WarningTok}[1]{\textcolor[rgb]{0.38,0.63,0.69}{\textbf{\textit{#1}}}}
\usepackage{longtable,booktabs,array}
\usepackage{calc} % for calculating minipage widths
% Correct order of tables after \paragraph or \subparagraph
\usepackage{etoolbox}
\makeatletter
\patchcmd\longtable{\par}{\if@noskipsec\mbox{}\fi\par}{}{}
\makeatother
% Allow footnotes in longtable head/foot
\IfFileExists{footnotehyper.sty}{\usepackage{footnotehyper}}{\usepackage{footnote}}
\makesavenoteenv{longtable}
\setlength{\emergencystretch}{3em} % prevent overfull lines
\providecommand{\tightlist}{%
  \setlength{\itemsep}{0pt}\setlength{\parskip}{0pt}}
\setcounter{secnumdepth}{5}
% Self-contained presentation rules embedded in every generated LaTeX file.
\usepackage{iftex}
\ifPDFTeX\else\PackageError{SSE-ASAS}{Compile with pdfLaTeX only}{}\fi
\usepackage[T1]{fontenc}
\usepackage[utf8]{inputenc}
\usepackage{lmodern}
\usepackage{microtype}
\usepackage{geometry}
\geometry{margin=0.93in,headheight=16pt}
\usepackage{xcolor} % [table] passed at top of file (option-clash fix)
\definecolor{ink}{HTML}{15283B}
\definecolor{brand}{HTML}{075B75}
\definecolor{brandlight}{HTML}{EAF5F8}
\definecolor{warm}{HTML}{A15C1D}
\definecolor{muted}{HTML}{536577}
\usepackage{fancyhdr}
\usepackage{titlesec}
\usepackage{enumitem}
\usepackage{booktabs}
\usepackage{longtable}
\usepackage{array}
\usepackage{tabularx}
\usepackage{tcolorbox}
\tcbuselibrary{breakable,skins}
\usepackage{listings}
\usepackage{fancyvrb}
\usepackage{hyperref}
\hypersetup{colorlinks=true,linkcolor=brand,urlcolor=brand,citecolor=brand}
\urlstyle{same}
\usepackage{lastpage}
\setlength{\emergencystretch}{3em}
\setlength{\parskip}{0.45em}
\setlength{\parindent}{0pt}
\setlist[itemize]{leftmargin=*,topsep=3pt,itemsep=2pt}
\setlist[enumerate]{leftmargin=*,topsep=3pt,itemsep=2pt}
\titleformat{\section}{\Large\bfseries\color{ink}}{\thesection}{0.75em}{}[\vspace{-2pt}\color{brand}\titlerule]
\titleformat{\subsection}{\large\bfseries\color{brand}}{\thesubsection}{0.7em}{}
\titleformat{\subsubsection}{\normalsize\bfseries\color{ink}}{\thesubsubsection}{0.7em}{}
\pagestyle{fancy}
\fancyhf{}
\fancyhead[L]{\small\sffamily\textcolor{brand}{SOFTWARE SECURITY ENGINEERING}}
\fancyhead[R]{\small\sffamily\textcolor{muted}{TECHNICAL SERIES}}
\fancyfoot[L]{\scriptsize\sffamily Version 1.0 \textbullet\ 14 October 2026}
\fancyfoot[R]{\scriptsize\sffamily Page \thepage\ of \pageref{LastPage}}
\renewcommand{\headrulewidth}{0.5pt}
\renewcommand{\footrulewidth}{0.2pt}
\setcounter{tocdepth}{2}
\setcounter{secnumdepth}{3}
\lstdefinestyle{sse}{basicstyle=\ttfamily\footnotesize,frame=single,framerule=0.4pt,rulecolor=\color{brand!40},backgroundcolor=\color{brandlight!45},numbers=left,numberstyle=\scriptsize\color{muted},numbersep=6pt,breaklines=true,breakatwhitespace=false,tabsize=2,showstringspaces=false,columns=fullflexible,keepspaces=true,captionpos=b,aboveskip=9pt,belowskip=9pt}
\lstset{style=sse}
\newtcolorbox{sourceconcept}[1][]{enhanced,breakable,colback=brandlight,colframe=brand!65!black,title={ORIGINAL CONCEPT},fonttitle=\bfseries\sffamily,#1}
\newtcolorbox{modernextension}[1][]{enhanced,breakable,colback=white,colframe=warm!85!black,title={MODERN EXTENSION},fonttitle=\bfseries\sffamily,#1}
\newtcolorbox{fieldwork}[1][]{enhanced,breakable,colback=white,colframe=ink!55,title={PRACTICAL APPLICATION},fonttitle=\bfseries\sffamily,#1}
\AtBeginDocument{\color{ink}}
\makeatletter
\renewcommand{\maketitle}{%
 \begin{titlepage}\thispagestyle{empty}
 \vspace*{1.5cm}{\color{brand}\rule{\textwidth}{3pt}}\par\vspace{1cm}
 {\sffamily\bfseries\color{muted}\large SOFTWARE SECURITY ENGINEERING\\ APPLICATION SECURITY ASSESSMENT SERIES}\par\vspace{1.05cm}
 {\sffamily\bfseries\color{ink}\fontsize{27}{34}\selectfont\@title\par}\vspace{0.8cm}
 {\sffamily\large\color{brand}\@author\par}\vfill
 \begin{tcolorbox}[colback=brandlight,colframe=brand!55,boxrule=0.5pt]
 \small\textbf{Editorial policy}\quad Original teaching synthesis with book-grounded concepts distinguished from contemporary practices. Modern platform behavior is version-sensitive; validate against linked vendor documentation before production use.
 \end{tcolorbox}
 {\color{brand}\rule{\textwidth}{2pt}}\par\vspace{0.3cm}
 {\sffamily\small Version 1.0 \quad|\quad October 2026 \quad|\quad pdfLaTeX}\end{titlepage}\setcounter{page}{1}}
\makeatother
\ifLuaTeX
  \usepackage{selnolig}  % disable illegal ligatures
\fi
\usepackage{bookmark}
\IfFileExists{xurl.sty}{\usepackage{xurl}}{} % add URL line breaks if available
\urlstyle{same}
\hypersetup{
  pdftitle={Hands-on Security Engineering Laboratories},
  pdfauthor={Independent Technical Study Edition},
  hidelinks,
  pdfcreator={LaTeX via pandoc}}

\title{Hands-on Security Engineering Laboratories}
\author{Independent Technical Study Edition}
\date{}

\begin{document}
\maketitle

{
\setcounter{tocdepth}{3}
\tableofcontents
}
\section{Laboratory handbook: rules, prerequisites, and
grading}\label{laboratory-handbook-rules-prerequisites-and-grading}

These exercises train safe reasoning and remediation. Use a local
throwaway repository, non-production fixtures, synthetic credentials,
and an isolated environment. Do not scan or test a system without
permission. Each lab creates an artifact and an evaluator can verify it
without guessing what the learner meant.

\textbf{Prerequisite order:} L01--L03 require Foundations. L04--L07
require Vulnerability Taxonomy and Manual Code Auditing. L08--L10
require Static Analysis Algorithms. L11--L14 require Secure Coding and
Assessment Procedures. \textbf{Source-derived concepts} come from the
three books; tooling, CI, CWE and ASVS integration are \textbf{modern
extensions}.

\section{L01: access control matrix
(foundational)}\label{l01-access-control-matrix-foundational}

\textbf{Goal:} differentiate authentication, permissions, and resource
ownership. \textbf{Time:} 45 minutes. \textbf{Inputs:} principals Alice
(tenant A, editor), Bob (tenant B, viewer), and Admin (cross-tenant
support with approval); objects invoices A1, B1, and policy
configuration.

\begin{enumerate}
\def\labelenumi{\arabic{enumi}.}
\tightlist
\item
  Build an access matrix for \texttt{read}, \texttt{update},
  \texttt{approve}, \texttt{export}.
\item
  Require that only a distinct approver approves an invoice.
\item
  Construct at least six denied access cases including cross-tenant
  export.
\item
  Explain whether an app gateway authenticating users establishes the
  required authorization.
\end{enumerate}

\textbf{Deliverable:} matrix, invariants, and test table. \textbf{Pass:}
all subjects/objects/rights covered; no implicit permit default;
ownership and role policies separated. \textbf{Source:} Bishop Chapters
2, 4, 8, 16. \textbf{Modern extension:} object-level API authorization.

\section{L02: information-flow and trust-boundary
analysis}\label{l02-information-flow-and-trust-boundary-analysis}

\textbf{Time:} 60 minutes. Map the flow from customer upload to object
storage, antivirus worker, message queue, search index, notification,
and logs. Label trust domains and where sensitive data could leak.
Identify three indirect flow paths not represented by an HTTP request.
Produce a diagram and a table of enforcement points, log redaction, and
detection controls.

\textbf{Pass:} each sensitive data store has a subject, required policy,
boundary, and verification test. \textbf{Source:} Bishop Chapters 5, 9,
17, 18; Dowd Chapter 2.

\section{L03: Clark--Wilson transaction
invariant}\label{l03-clarkwilson-transaction-invariant}

\textbf{Time:} 60 minutes. Model expense reimbursement as
\texttt{created\ -\textgreater{}\ submitted\ -\textgreater{}\ reviewed\ -\textgreater{}\ approved\ -\textgreater{}\ paid};
submitter must not approve, paid is terminal, and every approval is
attributable. Define state transitions, roles, and at least one
counterexample for a race between review and approval.

\textbf{Deliverable:} state table and six negative test cases.
\textbf{Pass:} all unauthorized transitions forbidden and at least one
atomicity mechanism specified. \textbf{Source:} Bishop Chapter 6.

\section{L04: classify ten flaws using distinct
vocabularies}\label{l04-classify-ten-flaws-using-distinct-vocabularies}

\textbf{Time:} 60 minutes. Catalog these observations: SQL string
concatenation; cross-tenant order lookup; shell command building; HTML
sink using untrusted markup; unchecked C multiplication; symlink race;
verbose error with token; public debug admin route; uncontrolled
decompression; duplicate webhook settlement. Assign lifecycle-origin
class, security property, likely CWE, and suggested verification
activity. \textbf{Do not} claim all have CVEs.

\textbf{Pass:} justifies every CWE against a real definition; identifies
when a reported problem is only a hypothesis. \textbf{Source:} Dowd
Chapters 1, 5--13, 17--18; Chess Chapters 5--12. \textbf{Modern
extension:} CWE and OWASP.

\section{L05: SQL injection root-cause
audit}\label{l05-sql-injection-root-cause-audit}

\textbf{Time:} 75 minutes. Work on the following local toy function.
Treat input as untrusted and the database as the relevant interpreter.

\begin{Shaded}
\begin{Highlighting}[]
\KeywordTok{def}\NormalTok{ find\_account(conn, name):}
\NormalTok{    q }\OperatorTok{=} \StringTok{"SELECT id FROM accounts WHERE name=\textquotesingle{}"} \OperatorTok{+}\NormalTok{ name }\OperatorTok{+} \StringTok{"\textquotesingle{}"}
    \ControlFlowTok{return}\NormalTok{ conn.execute(q).fetchall()}
\end{Highlighting}
\end{Shaded}

\begin{enumerate}
\def\labelenumi{\arabic{enumi}.}
\tightlist
\item
  Mark source, transformation, sink, and violated invariant.
\item
  Patch with bound values supported by your chosen database driver.
\item
  Test a normal value and one string containing a quote and SQL-like
  punctuation.
\item
  Document why HTML encoding is not an appropriate defense.
\end{enumerate}

\textbf{Pass:} normal inputs retain behavior, untrusted input cannot
alter SQL structure, and test uses only local synthetic data.
\textbf{Source:} Chess Chapter 5; Dowd Chapter 8.

\section{L06: cross-tenant ID access}\label{l06-cross-tenant-id-access}

\textbf{Time:} 90 minutes. Create an in-memory service containing
\texttt{order\_id}, \texttt{tenant\_id}, and \texttt{status}. Add
\texttt{get\_order(principal,\ order\_id)} that initially selects only
by ID. Write a failing negative test for tenant B's order and patch the
query to enforce the principal's tenant scope. Consider bulk list and
export functions as sibling paths.

\textbf{Deliverable:} policy, before/after diff, tests, and audit note.
\textbf{Pass:} a valid foreign ID is denied, no existence details leak,
and sibling paths are inspected. \textbf{Source:} Dowd Chapters 2--4.
\textbf{Modern extension:} API authorization testing.

\section{L07: integer bounds and native
parser}\label{l07-integer-bounds-and-native-parser}

\textbf{Time:} 90 minutes. Review \texttt{malloc(count\ *\ item\_size)}
followed by a loop that writes \texttt{count} records. Document
signedness, possible wraparound, source buffer length, allocation
failure, and loop termination. Implement a checked multiplication guard
and add tests at zero, maximum permitted length, and overflow boundary.

\begin{Shaded}
\begin{Highlighting}[]
\ControlFlowTok{if} \OperatorTok{(}\NormalTok{item\_size }\OperatorTok{==} \DecValTok{0} \OperatorTok{||}\NormalTok{ count }\OperatorTok{\textgreater{}}\NormalTok{ SIZE\_MAX }\OperatorTok{/}\NormalTok{ item\_size}\OperatorTok{)}
    \ControlFlowTok{return}\NormalTok{ ERROR}\OperatorTok{;}
\DataTypeTok{size\_t}\NormalTok{ total }\OperatorTok{=}\NormalTok{ count }\OperatorTok{*}\NormalTok{ item\_size}\OperatorTok{;}
\end{Highlighting}
\end{Shaded}

\textbf{Pass:} arithmetic is safe for input domain, source
readable-length invariant is explicitly checked, tests compile with
warnings and sanitizers where available. \textbf{Source:} Dowd Chapters
5--7; Chess Chapters 6--7.

\section{L08: CFG and forward may-taint
analysis}\label{l08-cfg-and-forward-may-taint-analysis}

\textbf{Time:} 90 minutes. Analyze:

\begin{Shaded}
\begin{Highlighting}[]
\NormalTok{1 x := trusted\_constant}
\NormalTok{2 if request.has(\textquotesingle{}q\textquotesingle{}):}
\NormalTok{3     x := request[\textquotesingle{}q\textquotesingle{}]}
\NormalTok{4 else:}
\NormalTok{5     x := \textquotesingle{}safe\textquotesingle{}}
\NormalTok{6 run\_query("SELECT ... " + x)}
\end{Highlighting}
\end{Shaded}

Draw blocks and edges, define taint facts and transfer functions, run
the worklist by hand to convergence, and identify reportable sink.
Describe one path-sensitive condition under which the path may not
execute. \textbf{Pass:} correct \texttt{IN/OUT} at all nodes and
accurate explanation of may-taint. \textbf{Source:} Chess Chapter 4.

\section{L09: build a minimal
analyzer}\label{l09-build-a-minimal-analyzer}

\textbf{Time:} 2 hours. Parse a deliberately tiny language with
statements \texttt{input\ x}, \texttt{copy\ y\ x},
\texttt{clean\_html\ x}, \texttt{sink\_html\ x}, \texttt{sink\_sql\ x}.
Track sets of tainted variables and flag a sink only when taint is
present and no context-valid cleansing occurred. Add a branch-join case.
Do not advertise the tool as a production-grade analyzer.

\textbf{Deliverable:} source, specification, five unit tests (at least
two negative), and one known limitation. \textbf{Pass:} deterministic
output with rule IDs and source locations. \textbf{Source:} Chess
Chapter 4; \textbf{modern extension:} SARIF ideas.

\section{L10: CodeQL workflow and SARIF
interpretation}\label{l10-codeql-workflow-and-sarif-interpretation}

\textbf{Time:} 2 hours. In an authorized small practice repo with a
CodeQL-supported language: enable the simplest supported setup, confirm
a successful analysis, list the actual language and file coverage,
inspect a harmless known-bad synthetic sample, and record the rule and
path. If private-repository licensing prevents scanning, demonstrate
locally on a suitable public practice repo or provide a written
configuration simulation without claiming an executed scan.

\textbf{Deliverable:} configuration snapshot, analysis run link or local
logs, alert explanation, limitations. \textbf{Pass:} distinguishes
configuration success from successful extraction; identifies run commit
and branch. \textbf{Modern extension:} CodeQL/GitHub Code Security.

\section{L11: HTML sink and safe DOM
APIs}\label{l11-html-sink-and-safe-dom-apis}

\textbf{Time:} 60 minutes. Build a browser-only practice page that
displays a supplied message via \texttt{innerHTML}. Replace with
\texttt{textContent} for plain text and test a value containing
markup-like characters. Document when HTML is intentionally required and
what additional controls become necessary. \textbf{Pass:} content
displayed literally without interpreted markup. \textbf{Source:} Chess
Chapter 9. \textbf{Modern extension:} current browser APIs/ASVS.

\section{L12: credential incident
tabletop}\label{l12-credential-incident-tabletop}

\textbf{Time:} 75 minutes. Use a fabricated string explicitly designated
\textbf{NOT A REAL TOKEN} in a temporary mock repo. Model secret
detection, incident triage, token rotation order, audit log review,
remediation of history, and a post-incident prevention control.
\textbf{No live credential should be committed, pasted, or printed.}

\textbf{Deliverable:} response timeline, roles, closure evidence.
\textbf{Pass:} distinguishes deleting text from revoking credential
validity. \textbf{Source:} Chess Chapter 11; Bishop assurance and
auditing; \textbf{modern extension:} GitHub Secret Protection.

\section{L13: manual review vs SAST coverage
challenge}\label{l13-manual-review-vs-sast-coverage-challenge}

\textbf{Time:} 2 hours. Provide two independent problems in a toy
service: (a) query construction from a request parameter and (b) an
approval policy that allows the submitter to approve their own
reimbursement. Predict which will be detected by configured static
queries and which requires invariant-aware review. Run the scanner if
entitled, compare prediction to output, then design both remediation
tests.

\textbf{Pass:} avoids inferring coverage from raw alert count; records
tested and untested hypotheses. \textbf{Source:} Dowd Chapters 2--4;
Chess Chapters 2--4. \textbf{Modern extension:} CodeQL results.

\section{L14: capstone assessment}\label{l14-capstone-assessment}

\textbf{Time:} 6--8 hours. Independently assess an intentionally
insecure local application with authentication, invoice API, file
upload, background job, and CI configuration. Follow the Assessment
Procedures guide end to end.

\textbf{Required artifacts:} authorized scope; threat model; entry-point
inventory; scanning coverage; manual audit notes; at least three
distinct root-cause findings or justified no-findings with coverage
limits; one patch and regression test; risk/disposition ledger;
signed-off residual-risk statement; an executive one-page summary.

\textbf{Pass:} every claimed finding has immutable source location,
trace, observed evidence or clearly labeled hypothesis, impact, precise
weakness category, owner, and fix verification. No destructive proof of
concept is needed.

\section{Assessment rubric and answer
checkpoints}\label{assessment-rubric-and-answer-checkpoints}

\begin{longtable}[]{@{}
  >{\raggedright\arraybackslash}p{(\columnwidth - 4\tabcolsep) * \real{0.3333}}
  >{\raggedleft\arraybackslash}p{(\columnwidth - 4\tabcolsep) * \real{0.3333}}
  >{\raggedright\arraybackslash}p{(\columnwidth - 4\tabcolsep) * \real{0.3333}}@{}}
\toprule\noalign{}
\begin{minipage}[b]{\linewidth}\raggedright
Criterion
\end{minipage} & \begin{minipage}[b]{\linewidth}\raggedleft
Weight
\end{minipage} & \begin{minipage}[b]{\linewidth}\raggedright
Observable evidence
\end{minipage} \\
\midrule\noalign{}
\endhead
\bottomrule\noalign{}
\endlastfoot
Policy/invariant clarity & 20\% & Specific denied behavior and protected
asset \\
Trace correctness & 20\% & Correct source, path, sink, authorization
context \\
Remediation quality & 20\% & Root cause removed with preserved
functionality \\
Regression evidence & 20\% & Test would fail pre-fix and pass
post-fix \\
Scope and ethical handling & 10\% & Authorized, synthetic data, clear
limits \\
Reporting quality & 10\% & Reproducible, versioned artifacts and
mapping \\
\end{longtable}

\textbf{Answer checkpoints (not turnkey solutions):} L01: cross-tenant
access must fail even for a privileged UI user without scope; L05: the
SQL syntax must remain constant as input changes; L07: guard must
execute \emph{before} multiplication; L08: join merges the tainted
branch with the untainted branch; L09: a HTML-specific cleaner must not
cleanse an SQL sink; L13: missing business-logic authorization may not
produce a scanner alert. Grade the reasoning, not just the number of
findings.

\section{Lab evidence manifest}\label{lab-evidence-manifest}

For each lab, store \texttt{lab\_id}, training repo SHA, environment,
tool versions, command or proof-of-review, results, tested invariant,
associated CWE/ASVS if appropriate, reviewer, residual limitations, and
completion status. Keep any screenshots redacted. When machine execution
is not available, label the deliverable as a \textbf{design exercise}
rather than falsifying a scan run.

\section{Annex A: reusable lab fixtures and evidence
templates}\label{annex-a-reusable-lab-fixtures-and-evidence-templates}

\subsection{Lab workspace convention}\label{lab-workspace-convention}

Create one local directory per lab. A suggested layout is
\texttt{README.md} for authorized scope and setup, \texttt{src/} for
deliberately small sample implementations, \texttt{tests/} for negative
and positive tests, \texttt{evidence/} for redacted tool results, and
\texttt{review.md} for policy, root cause, fix and limitations. Check
only synthetic credentials and synthetic user data into fixtures. Build
scripts should not download untrusted executables or call outside
services unless explicitly approved.

For each lab, start by checking the \emph{expected vulnerable state}
using a test that describes the undesired behavior. Apply the smallest
reasonable design fix, rerun the same test, and add adjacent cases. This
before/after procedure prevents a learner from crediting a patch that
never exercised the flaw.

\subsection{L06 fixture: object
authorization}\label{l06-fixture-object-authorization}

\begin{Shaded}
\begin{Highlighting}[]
\OtherTok{[}
  \FunctionTok{\{}\DataTypeTok{"id"}\FunctionTok{:} \DecValTok{101}\FunctionTok{,} \DataTypeTok{"tenant"}\FunctionTok{:} \StringTok{"A"}\FunctionTok{,} \DataTypeTok{"owner"}\FunctionTok{:} \StringTok{"alice"}\FunctionTok{,} \DataTypeTok{"amount"}\FunctionTok{:} \DecValTok{25}\FunctionTok{\}}\OtherTok{,}
  \FunctionTok{\{}\DataTypeTok{"id"}\FunctionTok{:} \DecValTok{102}\FunctionTok{,} \DataTypeTok{"tenant"}\FunctionTok{:} \StringTok{"B"}\FunctionTok{,} \DataTypeTok{"owner"}\FunctionTok{:} \StringTok{"bob"}\FunctionTok{,} \DataTypeTok{"amount"}\FunctionTok{:} \DecValTok{50}\FunctionTok{\}}
\OtherTok{]}
\end{Highlighting}
\end{Shaded}

\textbf{Positive:} Alice can read invoice 101 while holding the correct
action permission. \textbf{Negative 1:} Alice cannot read 102 even
though the ID exists. \textbf{Negative 2:} a principal without the read
permission cannot read 101. \textbf{Negative 3:} the API does not return
a list whose contents include B's record. \textbf{Negative 4:} an export
worker cannot fetch 102 on Alice's behalf. The same root cause may
appear in four code paths; each requires an explicit test or a shared
control with proven reuse.

\subsection{L07 fixture: numerical
boundaries}\label{l07-fixture-numerical-boundaries}

Test the allocation expression with zero count, zero item size, one
record, maximum permitted count, \texttt{SIZE\_MAX\ /\ item\_size}, and
one more than that representable bound when expressible in the input
type. Assert rejection before multiplication for non-representable
products. Also verify that input buffer length is large enough and that
zero-sized allocations have deliberate handling. Run compiler warnings
and sanitizers as supported by your compiler/architecture, documenting
unavailable tools.

\subsection{L08 expected facts}\label{l08-expected-facts}

At the entry to a simple may-taint analysis, the fact set is empty. The
request-derived assignment adds a taint fact. At a branch merge, join
includes it even if another branch assigns a constant. Only a valid
sanitizer modeled on all reachable paths can remove the property for the
specific sink. If the learner asserts that reaching the sink proves
exploitability, require them to revisit path feasibility and deployment
preconditions.

\subsection{L09 analyzer test
inventory}\label{l09-analyzer-test-inventory}

\begin{longtable}[]{@{}
  >{\raggedright\arraybackslash}p{(\columnwidth - 4\tabcolsep) * \real{0.3333}}
  >{\raggedright\arraybackslash}p{(\columnwidth - 4\tabcolsep) * \real{0.3333}}
  >{\raggedright\arraybackslash}p{(\columnwidth - 4\tabcolsep) * \real{0.3333}}@{}}
\toprule\noalign{}
\begin{minipage}[b]{\linewidth}\raggedright
Test
\end{minipage} & \begin{minipage}[b]{\linewidth}\raggedright
Input program
\end{minipage} & \begin{minipage}[b]{\linewidth}\raggedright
Expected assessment
\end{minipage} \\
\midrule\noalign{}
\endhead
\bottomrule\noalign{}
\endlastfoot
S01 & \texttt{input\ x;\ sink\_sql\ x} & SQL-flow finding \\
S02 & \texttt{input\ x;\ clean\_html\ x;\ sink\_sql\ x} & SQL-flow
finding remains \\
S03 & \texttt{input\ x;\ clean\_html\ x;\ sink\_html\ x} & No HTML-flow
finding under model \\
S04 & \texttt{copy\ y\ x;\ sink\_html\ y} after \texttt{input\ x} &
HTML-flow finding \\
S05 & \texttt{x=constant;\ sink\_sql\ x} & No taint finding \\
S06 & Two-branch merge: one branch tainted & May-flow finding \\
\end{longtable}

This toy grammar is deliberately limited and may not model aliasing or
real control flow; learners must state limits prominently. They may use
a parser library or a hand-written tokenizer, but the model's semantics
must be documented.

\section{Annex B: capstone review roles and
artifacts}\label{annex-b-capstone-review-roles-and-artifacts}

\textbf{Learner / assessor:} proposes hypotheses and creates evidence.
\textbf{Engineering reviewer:} checks that code and regression tests
establish claimed fixes. \textbf{Security reviewer:} evaluates trust
assumptions, weakness mapping, exploitability conditions, and gaps.
\textbf{Risk approver:} accepts remaining business exposure. One person
may hold multiple training roles, but the roles' independent questions
must still be answered in writing.

Deliver a fixed bundle: \texttt{scope}, \texttt{asset\ inventory},
\texttt{policy\ and\ threat\ model}, \texttt{audit\ cycles},
\texttt{scanner\ configuration\ +\ coverage}, \texttt{findings\ ledger},
\texttt{patch\ and\ tests}, \texttt{risk\ decision}, and
\texttt{one-page\ summary}. The bundle should be understandable even if
no scan could run because licensing or networking was unavailable;
explain the limitation and use a simulation only when explicitly labeled
as such.

\textbf{Oral defense questions:} Which security property failed? Why was
the path reachable? What evidence would falsify your finding? Why is the
chosen fix more comprehensive than the local patch? Which sibling
components remain unassessed? What did the tool \emph{actually} scan?
Can an exception be justified without hiding the weakness?

\section{Annex C: competency remediation
plan}\label{annex-c-competency-remediation-plan}

If a learner fails \textbf{policy modeling}, repeat L01--L03 with
different assets and actors. If they fail \textbf{mechanism
classification}, do L04 with fewer cases and official CWE definitions.
If they fail \textbf{manual trace correctness}, perform L05--L07 with
peer review and a prescribed call-chain worksheet. If they fail
\textbf{static-analysis reasoning}, rerun L08--L09 with explicit lattice
and transfer tables. If they fail \textbf{evidence integrity}, require
resubmission with exact code revision, environment, and verified logs;
do not award points for screenshots without reproducibility.

\section{Source and modern extension
map}\label{source-and-modern-extension-map}

\textbf{Book-grounded:} Bishop Chapters 2--9, 17, 19--25; Dowd Chapters
1--4, 5--13, 17--18; Chess \& West Chapters 2--4, 5--12, 13--14
(source-code analysis exercises).

\textbf{Modern extension:} \url{https://cwe.mitre.org/};
\url{https://owasp.org/projects/asvs};
\url{https://docs.github.com/en/code-security/concepts/code-scanning/codeql/codeql-cli};
\url{https://docs.github.com/en/code-security/concepts/secret-security/about-alerts};
\url{https://csrc.nist.gov/pubs/sp/800/218/final}.

\end{document}
